% Cover body for the palette pitch. Colour roles are set by the wrapper: ColTorsion (t), ColUmbrella (tu on the
% band), ColStar (b, the starred generator, between twin sheets), ColDot (versions), ColBand/ColBandLine (core
% sheet), ColWall (cladding), ColRetained (the packet), plus the line styles tpath, upath, lift, blift, vloopt, vloopb.
\begin{tikzpicture}[plate,head/.style={font=\fontsize{9.5}{11}\selectfont\bfseries,inner sep=0pt},rng/.style={font=\fontsize{8}{10}\selectfont,inner sep=0pt}]
\path[use as bounding box] (0,0) rectangle (15,14.4);
\pgfmathsetmacro{\Ro}{2.3}\pgfmathsetmacro{\Ri}{1.62}\pgfmathsetmacro{\E}{.42}
% ------------------------------------------------ the slab C with the hole the solid leaves
\begin{scope}[shift={(0,.9)}]
\path[cut] (.5,.9)--(12.1,.9)--(14.5,3.1)--(2.9,3.1)--cycle; \draw[heavy] (.5,.9)--(12.1,.9)--(14.5,3.1)--(2.9,3.1)--cycle;
\node[tag,anchor=north west] at (3.05,3.0) {$\conf$};
\path[fill=white,even odd rule] (7.6,2.0) ellipse[x radius=\Ro,y radius={\E*\Ro}] (7.6,2.0) ellipse[x radius=\Ri,y radius={\E*\Ri}];
\draw[fine] (7.6,2.0) ellipse[x radius=\Ro,y radius={\E*\Ro}]; \draw[fine] (7.6,2.0) ellipse[x radius=\Ri,y radius={\E*\Ri}];
% ------------------------------------------------ one configuration, a point of C
\fill[PlateInk] (2.7,2.35) circle[radius=1.4pt];
\begin{scope}\molsolid\mollabelsfalse\pic[scale=.72] at (2.0,5.0) {mol3d-mla-ref};\end{scope}
\draw[leader] (2.25,4.1)--(2.7,2.45);
\node[head,anchor=west] at (.5,7.5) {Configurations and States}; \node[rng,anchor=north west,text width=4.2cm] at (.5,7.28) {\S\S\,1--3\quad a configuration is a point of $\conf$};
% ------------------------------------------------ the feasible solid F, its band, the two generator paths, the versions, a packet
\begin{scope}[shift={(7.6,5.3)}]\tikzset{tube/Ro=2.3,tube/Ri=1.62,tube/H=.55,tube/E=.42,tube/a1=-62,tube/a2=8}
  \draw[hidden] (-\Ro,-3.3)--(-\Ro,-.55) (\Ro,-3.3)--(\Ro,-.55);
  \pic (ft) at (0,0) {thicktube};
  \draw[tpath] plot[domain=0:115,samples=30,variable=\a] ({1.96*sin(\a)},{.55-.42*1.96*cos(\a)});
  \draw[upath] plot[domain=0:1,samples=30,variable=\s] ({1.96*sin(-60*\s)},{.55-2*.55*\s-.42*1.96*cos(60*\s)});
  \path[fill=ColRetained!55] (ft-top0) circle[radius=.3]; \draw[fine] (ft-top0) circle[radius=.3];
  \foreach \p in {top0,top120,top240,bot300} {\hatomat{ft-\p}{.085}{ColDot}}
  \foreach \p in {bot60,bot180} {\draw[ghost,fill=ColDot!45] (ft-\p) circle[radius=.085];}
  \node[tag,anchor=west] at (2.5,-1.15) {$\mathcal F$};
\end{scope}
\node[head,anchor=west] at (10.7,5.2) {Fragments and}; \node[head,anchor=west] at (10.7,4.8) {Feasible Regimes}; \node[rng,anchor=north west,text width=3.9cm] at (10.7,4.58) {\S\S\,4--6\quad $\mathcal F$ carved out of $\conf$, its versions on the band};
\node[head,anchor=west] at (10.7,7.4) {Symmetry and}; \node[head,anchor=west] at (10.7,7.0) {Localization}; \node[rng,anchor=north west,text width=3.9cm] at (10.7,6.78) {\S\S\,7--9\quad species on the versions, a packet on one of them};
\draw[leader] (10.6,7.2)--(9.55,6.45);
% ------------------------------------------------ the quotient patch V with the two loops, the six sheets over it with the two lifts
\draw[fine,dash pattern=on 1.6pt off 1.3pt] (7.6,5.03)--(7.9,8.52);
\node[sub,anchor=west,inner sep=1pt] at (7.85,7.3) {$\mathcal F\to\mathcal F/G_{12}$};
\begin{scope}[shift={(7.6,8.3)}]
  \path[fill=white] (-1.2,0)--(1.2,0)--(1.8,.45)--(-.6,.45)--cycle; \draw[heavy] (-1.2,0)--(1.2,0)--(1.8,.45)--(-.6,.45)--cycle;
  \coordinate (bx) at (.3,.22); \fill[PlateInk] (bx) circle[radius=1.2pt];
  \draw[vloopt] (bx) .. controls (-.9,.6) and (-.9,-.15) .. (bx);
  \draw[vloopb] (bx) .. controls (1.5,-.15) and (1.5,.6) .. (bx);
  \node[sub,anchor=west,inner sep=1pt] at (1.9,.22) {$V$};
  \foreach \k in {0,1,2} {\pgfmathsetmacro{\ya}{1.15+1.2*\k}\pgfmathsetmacro{\yb}{\ya+.32}
    \path[fill=white] (-1.2,\ya)--(1.2,\ya)--(1.8,\ya+.45)--(-.6,\ya+.45)--cycle; \draw[fine] (-1.2,\ya)--(1.2,\ya)--(1.8,\ya+.45)--(-.6,\ya+.45)--cycle;
    \path[fill=white] (-1.2,\yb)--(1.2,\yb)--(1.8,\yb+.45)--(-.6,\yb+.45)--cycle; \draw[fine] (-1.2,\yb)--(1.2,\yb)--(1.8,\yb+.45)--(-.6,\yb+.45)--cycle;
    \coordinate (pa\k) at (.3,{\ya+.22}); \coordinate (pb\k) at (.3,{\yb+.22});}
  \foreach \k in {0,1,2} {\fill[PlateInk] (pa\k) circle[radius=1pt]; \fill[PlateInk] (pb\k) circle[radius=1pt];}
  \draw[lift] (pa0) .. controls (-.9,{1.15+.22+.4}) and (-.9,{1.15+1.2+.22-.4}) .. (pa1);
  \draw[blift] (pa1) .. controls (1.4,{1.15+1.2+.22+.05}) and (1.4,{1.15+1.2+.32+.22-.05}) .. (pb1);
  \draw[blift] (pa0) .. controls (1.4,{1.15+.22+.05}) and (1.4,{1.15+.32+.22-.05}) .. (pb0);
  \draw[lift] (pb0) .. controls (-1.0,{1.15+.32+.22+.8}) and (-1.0,{1.15+2.4+.32+.22-.8}) .. (pb2);
  \hatomat{pb1}{.075}{ColDot} \hatomat{pb2}{.075}{ColDot}
\end{scope}
\node[head,anchor=west] at (10.7,12.1) {Representations and}; \node[head,anchor=west] at (10.7,11.7) {Global Structure}; \node[rng,anchor=north west,text width=3.9cm] at (10.7,11.48) {\S\S\,10--12\quad charts, sheets and modes; the lifts of the two loops end on different sheets};
\draw[leader] (10.6,11.9)--(9.5,11.6);
\end{scope}
% ------------------------------------------------ the role strip
\node[sub,anchor=west] at (.5,1.3) {\paletteName};
\begin{scope}[shift={(.5,.55)}]
  \draw[tpath,-] (0,0)--(.7,0); \node[sub,anchor=west] at (.8,0) {$t$};
  \draw[upath,-] (1.5,0)--(2.2,0); \node[sub,anchor=west] at (2.3,0) {$tu$ on the band};
  \draw[blift] (4.5,0)--(5.2,0); \node[sub,anchor=west] at (5.3,0) {$b$, twin sheets};
  \hatom{7.6}{0}{.085}{ColDot} \node[sub,anchor=west] at (7.8,0) {versions};
  \path[fill=ColBand] (9.3,-.13) rectangle (9.8,.13); \draw[fine] (9.3,-.13) rectangle (9.8,.13); \node[sub,anchor=west] at (9.9,0) {core sheet};
  \path[fill=ColRetained!55] (11.7,-.13) rectangle (12.2,.13); \draw[fine] (11.7,-.13) rectangle (12.2,.13); \node[sub,anchor=west] at (12.3,0) {packet};
  \hatom{13.6}{0}{.11}{ColNitrogen} \node[sub,anchor=west] at (13.8,0) {N};
\end{scope}
\end{tikzpicture}
